\documentclass[../DoAn.tex]{subfiles}
\begin{document}

Chương trước đã nêu bối cảnh, mục tiêu và phạm vi của đề tài. Chương này khảo sát hiện trạng công tác thi đua khen thưởng tại Học viện Khoa học Quân sự (HVKHQS) và rút ra các yêu cầu chức năng, phi chức năng của phần mềm.
Từ thực trạng và những bất cập của cách quản lý thủ công, hệ thống được mô hình hoá qua cấu trúc tác nhân cùng biểu đồ use case tổng quát và phân rã; những use case tiêu biểu được đặc tả chi tiết.
Phần cuối tổng hợp các yêu cầu phi chức năng về hiệu năng, bảo mật và tính nhất quán dữ liệu, làm tiền đề cho thiết kế ở các chương sau.

\section{Khảo sát hiện trạng}
\label{section:2.1}

Trước khi triển khai phần mềm chuyên dụng, công tác thi đua khen thưởng tại HVKHQS phụ thuộc vào ba nguồn dữ liệu: hồ sơ giấy lưu tại Phòng Chính trị, các tệp chia sẻ qua mạng nội bộ và sổ theo dõi cùng các tệp lưu riêng của từng đơn vị.
Qua trao đổi với cán bộ phụ trách, các đặc điểm chính như sau.

Cấp đề xuất là chỉ huy các đơn vị.
Nhóm này nắm danh sách quân nhân và lập đề nghị theo định kỳ hoặc các đợt đặc biệt.
Cấp xét duyệt là Phòng Chính trị Học viện, kiểm tra điều kiện rồi lập tờ trình gửi Ban Giám đốc.
Cấp ký quyết định tùy theo thẩm quyền của từng danh hiệu: Giám đốc Học viện ký các danh hiệu trong thẩm quyền, còn những danh hiệu vượt thẩm quyền do cấp cao hơn ký, như Tổng cục, Bộ Quốc phòng, Thủ tướng Chính phủ hay Chủ tịch nước.
Quân nhân nhận thông tin khen thưởng chủ yếu theo một chiều qua văn bản hoặc bảng tin đơn vị, chưa có kênh tra cứu chủ động.

Mỗi quân nhân thường có nhiều quyết định khen thưởng từ các loại danh hiệu khác nhau: chuỗi danh hiệu hằng năm tích luỹ toàn thời gian phục vụ, Huy chương Chiến sĩ vẻ vang (HCCSVV) ở các mốc 10, 15, 20 năm, Huân chương Bảo vệ Tổ quốc (HCBVTQ), Kỷ niệm chương, Huy chương Quân kỳ quyết thắng và khen thưởng đột xuất.
Cán bộ phục vụ trên 25 năm có thể tích luỹ nhiều quyết định khen thưởng.
Ở quy mô vài nghìn cán bộ, học viên và sĩ quan, mỗi người lại tích lũy nhiều bản ghi khen thưởng qua các năm, nên tổng dữ liệu lịch sử là một khối lượng lớn.

Từ kết quả khảo sát, rút ra bốn nhóm vấn đề chính:

\begin{enumerate}
    \item Việc tra cứu lịch sử khen thưởng cá nhân còn mất nhiều thời gian do cán bộ phải mở và đối chiếu nhiều tài liệu và các tệp khác nhau.
    \item Đối chiếu chuỗi danh hiệu hằng năm dễ sai sót vì mỗi danh hiệu chỉ được xét trên một khoảng thời gian gần nhất chứ không phải toàn bộ quá trình. Khoảng xét lại lùi dần theo năm, có danh hiệu nhận lặp lại sau mỗi chu kỳ còn BKTTCP cá nhân chỉ nhận một lần, nên khi tính tay rất dễ nhầm.
    \item Chưa có cơ chế phát hiện trùng lặp khi nhiều cán bộ cùng đề nghị một quân nhân, hoặc khi đề nghị một danh hiệu mà quân nhân đã được nhận.
    \item Quy trình thủ công chưa có cơ chế lưu vết chỉnh sửa, gây khó truy vết khi phát hiện sai sót.
\end{enumerate}

Các đặc thù nghiệp vụ trên đòi hỏi một hệ thống có thể lưu trữ tập trung và liên kết các loại bản ghi, tự động kiểm tra điều kiện theo bộ quy tắc định sẵn, phân quyền theo cơ cấu tổ chức, và ghi nhật ký cho mọi thao tác làm thay đổi dữ liệu nghiệp vụ.

\section{Tổng quan chức năng}
\label{section:2.2}

Phần này trình bày cấu trúc tác nhân và các nhóm chức năng chính của hệ thống thông qua biểu đồ use case tổng quát và các biểu đồ phân rã chức năng.

\subsection{Biểu đồ use case tổng quát}
\label{subsection:2.2.1}

Hệ thống có bốn nhóm tác nhân tương ứng với bốn cấp phân quyền nội bộ:

\begin{itemize}
    \item \textbf{Quản trị viên (SuperAdmin)} đại diện cho bộ phận kỹ thuật và công nghệ thông tin. Nhóm này không tham gia trực tiếp vào nghiệp vụ xét khen thưởng mà tập trung vào quản trị hạ tầng, sao lưu và khôi phục dữ liệu, theo dõi nhật ký hệ thống và quản trị các tài khoản.
    \item \textbf{Cán bộ Phòng Chính trị (Admin)} là nhân sự thuộc Phòng Chính trị Học viện. Nhóm này có toàn quyền nghiệp vụ: phê duyệt đề xuất, quản lý quân nhân, đơn vị, chức vụ, quản lý các tệp quyết định khen thưởng, thêm số quyết định cho đề xuất khi phê duyệt và xuất danh sách khen thưởng.
    \item \textbf{Chỉ huy đơn vị (Manager)} là cán bộ chỉ huy cấp đơn vị, phụ trách lập đề xuất khen thưởng cho phạm vi đơn vị họ quản lý và theo dõi trạng thái phê duyệt. Chỉ huy đơn vị có thể thao tác trên dữ liệu quân nhân thuộc đơn vị nhưng không thể chuyển quân nhân giữa các đơn vị, vì chức năng này chỉ thuộc Cán bộ Phòng Chính trị.
    \item \textbf{Người dùng (User)} là cán bộ, học viên hoặc sĩ quan sử dụng hệ thống để tra cứu hồ sơ khen thưởng cá nhân, xem các mốc danh hiệu hằng năm đủ điều kiện đề nghị theo định kỳ và nhận thông báo về kết quả phê duyệt.
\end{itemize}

Hệ thống cung cấp các nhóm chức năng bao trùm quản lý hồ sơ quân nhân, đơn vị và tài khoản, khen thưởng cùng đề xuất và phê duyệt, kiểm tra điều kiện và gợi ý, quản lý quyết định, thông báo thời gian thực, báo cáo thống kê, nhật ký và sao lưu. Nhóm Quản lý khen thưởng bao gồm nhiều loại danh hiệu; bảy loại đề xuất tương ứng được trình bày chi tiết ở Mục~\ref{subsection:2.2.4}.
Mối quan hệ giữa từng tác nhân và các nhóm chức năng được thể hiện ở Hình \ref{fig:usecase-tongquat}.

\begin{figure}[H]
    \centering
    \includegraphics[width=\textwidth,height=0.9\textheight,keepaspectratio]{uc-tong-quan}
    \caption{Biểu đồ use case tổng quát hệ thống PM QLKT}
    \label{fig:usecase-tongquat}
\end{figure}

\subsection{Biểu đồ use case phân rã chức năng quản lý quân nhân và lịch sử chức vụ}
\label{subsection:2.2.2}

Quân nhân là thực thể trung tâm, liên kết tới hầu hết các bảng nghiệp vụ (danh hiệu hằng năm, HCCSVV, HCBVTQ, thành tích khoa học, khen thưởng đột xuất, đề xuất).
Phân rã chức năng gồm các use case con thể hiện ở Hình~\ref{fig:usecase-quannhan}:

\begin{itemize}
    \item \textbf{Tạo mới quân nhân}: chỉ Cán bộ Phòng Chính trị có quyền. Quân nhân được khởi tạo cùng lúc với tài khoản: cán bộ nhập tên đăng nhập, chọn vai trò, đơn vị và chức vụ; hệ thống kiểm tra trùng tên đăng nhập rồi tạo bản ghi quân nhân cùng tài khoản; cả hai hoặc cùng được tạo, hoặc cùng không khi có lỗi. Họ tên đầy đủ, ngày sinh, ngày nhập ngũ và các thông tin khác được bổ sung sau qua chức năng cập nhật.
    \item \textbf{Cập nhật thông tin}: Cán bộ Phòng Chính trị sửa được toàn bộ. Chỉ huy đơn vị sửa được thông tin chung của quân nhân thuộc đơn vị họ nhưng không được đổi đơn vị và chức vụ. Bản thân quân nhân chỉ tự cập nhật được một số thông tin cá nhân cơ bản: họ tên, ngày sinh, số điện thoại, quê quán, trú quán và chỗ ở hiện nay.
    \item \textbf{Chuyển đơn vị}: chỉ Cán bộ Phòng Chính trị. Khi chuyển, hệ thống cập nhật số lượng quân nhân của đơn vị cũ (giảm 1) và đơn vị mới (tăng 1); hoặc cả hai thay đổi cùng thành công, hoặc cùng được hoàn tác khi có lỗi. Toàn bộ lịch sử khen thưởng của quân nhân được giữ nguyên, không tạo bản ghi mới hay xoá bản ghi cũ.
    \item \textbf{Quản lý lịch sử chức vụ}: Cán bộ Phòng Chính trị nhập danh sách các giai đoạn quân nhân giữ chức vụ với hệ số tương ứng, gồm cả giai đoạn hiện tại đang đảm nhiệm. Chỉ huy đơn vị chỉ được bổ sung các giai đoạn chức vụ đã kết thúc và bắt buộc có ngày kết thúc, không được thêm giai đoạn hiện tại đang mở vì đổi chức vụ hiện tại thuộc quyền Cán bộ Phòng Chính trị. Dữ liệu này là điều kiện xét khen thưởng HCBVTQ.
    \item \textbf{Thêm/sửa/xoá thủ công qua biểu mẫu}: dữ liệu quân nhân được quản lý trực tiếp qua biểu mẫu trên giao diện.
\end{itemize}

\begin{figure}[H]
    \centering
    \includegraphics[width=\textwidth,height=0.9\textheight,keepaspectratio]{uc-quan-nhan}
    \caption{Biểu đồ use case phân rã chức năng quản lý quân nhân}
    \label{fig:usecase-quannhan}
\end{figure}

\subsection{Biểu đồ use case phân rã chức năng quản lý đơn vị và chức vụ}
\label{subsection:2.2.3}

Cấu trúc tổ chức của Học viện được mô hình hoá thành cơ cấu hai cấp: cơ quan đơn vị ở cấp trên và đơn vị trực thuộc ở cấp dưới.
Mỗi đơn vị trực thuộc luôn thuộc đúng một cơ quan đơn vị cha; ngược lại, một cơ quan đơn vị có thể có nhiều đơn vị trực thuộc hoặc chưa có đơn vị nào bên dưới.
Mỗi quân nhân luôn gắn với một cơ quan đơn vị; nếu công tác ở một đơn vị trực thuộc thì gắn thêm với đơn vị đó.

Nhóm chức năng quản lý đơn vị bao gồm các use case con: tạo mới cơ quan đơn vị, tạo mới đơn vị trực thuộc thuộc một cơ quan đơn vị cụ thể, cập nhật thông tin, xoá (chỉ thực hiện được khi không còn quân nhân nào thuộc đơn vị) và xem cơ cấu tổ chức dưới dạng cây phân cấp trên giao diện.
Danh mục chức vụ được quản lý kèm hệ số chức vụ nhận giá trị từ 0 đến 1{,}0 với một chữ số thập phân tuỳ chức danh; mỗi chức vụ gắn với cơ quan đơn vị hoặc đơn vị trực thuộc mà nó thuộc về.
Hệ số chức vụ là dữ liệu đầu vào trực tiếp cho việc tính thời gian giữ chức vụ phục vụ xét HCBVTQ theo nhóm hệ số 0{,}7, 0{,}8 và 0{,}9 đến 1{,}0.
Các use case của nhóm chức năng quản lý đơn vị và chức vụ được thể hiện ở Hình~\ref{fig:usecase-donvi}.

\begin{figure}[H]
    \centering
    \includegraphics[width=\textwidth,height=0.9\textheight,keepaspectratio]{uc-don-vi}
    \caption{Biểu đồ use case phân rã chức năng quản lý đơn vị và chức vụ}
    \label{fig:usecase-donvi}
\end{figure}

\subsection{Biểu đồ use case phân rã chức năng quản lý khen thưởng}
\label{subsection:2.2.4}

Đề xuất khen thưởng là điểm khởi đầu của mọi quy trình ghi nhận thành tích.
Hệ thống hỗ trợ bảy loại đề xuất riêng biệt, mỗi loại có dữ liệu đầu vào và quy tắc xét điều kiện khác nhau:

\begin{enumerate}
    \item \textbf{Cá nhân hằng năm}: gồm hai danh hiệu thi đua cấp cơ sở là Chiến sĩ tiên tiến (CSTT) và Chiến sĩ thi đua cơ sở (CSTĐCS), cùng ba mức cao hơn xét theo chuỗi nhiều năm là BKBQP, CSTĐTQ và BKTTCP.
    \item \textbf{Đơn vị hằng năm}: gồm hai danh hiệu thi đua cấp cơ sở là Đơn vị tiên tiến và Đơn vị quyết thắng (ĐVQT), cùng hai mức cao hơn xét theo chuỗi nhiều năm là BKBQP và BKTTCP.
    \item \textbf{HCCSVV}: Huy chương Chiến sĩ vẻ vang xét theo niên hạn phục vụ tích luỹ tới năm xét, chia ba hạng tăng dần là hạng Ba, hạng Nhì và hạng Nhất ứng với 10, 15 và 20 năm. Các hạng xét tuần tự, phải có hạng dưới trước khi đề nghị hạng trên.
    \item \textbf{HCBVTQ}: Huân chương Bảo vệ Tổ quốc xét theo thời gian giữ chức vụ thuộc các nhóm hệ số, cũng chia ba hạng tăng dần. Hạng Ba cộng thời gian giữ chức ở nhóm hệ số từ 0{,}7 trở lên, hạng Nhì từ 0{,}8 trở lên và hạng Nhất chỉ tính nhóm 0{,}9 đến 1{,}0. Thời gian giữ chức tối thiểu là 10 năm với nam, nữ được giảm một phần ba, và các hạng cũng xét tuần tự từ thấp lên cao.
    \item \textbf{Kỷ niệm chương}: hiện có Kỷ niệm chương vì sự nghiệp xây dựng QĐNDVN, mỗi quân nhân chỉ được trao một lần khi đủ số năm phục vụ; ngưỡng yêu cầu là 25 năm với nam và 20 năm với nữ.
    \item \textbf{Quân kỳ quyết thắng}: Huy chương Quân kỳ quyết thắng cho cán bộ phục vụ đủ 25 năm.
    \item \textbf{Thành tích khoa học}: Khen thưởng đề tài khoa học hoặc sáng kiến khoa học.
\end{enumerate}

Cả bảy loại có cấu trúc dữ liệu và quy tắc xét riêng nhưng dùng chung một quy trình đề xuất và phê duyệt. Chi tiết cách tổ chức được trình bày ở Chương~\ref{chapter:Experiment}.

Trên giao diện, quy trình tạo đề xuất gồm năm bước nối tiếp: chọn loại đề xuất; chọn năm hoặc tháng xét cùng quân nhân hoặc đơn vị; thiết lập danh hiệu cho từng đối tượng; tải lên tệp đính kèm; cuối cùng xem lại và gửi đề xuất.
Bước hai cho phép chọn hàng loạt qua các bộ lọc theo đơn vị, năm sinh, giới tính và mốc thời gian phục vụ.
Toàn bộ use case của nhóm chức năng quản lý khen thưởng được thể hiện ở Hình~\ref{fig:usecase-dexuat}.

\begin{figure}[H]
    \centering
    \includegraphics[width=\textwidth,height=0.9\textheight,keepaspectratio]{uc-khen-thuong}
    \caption{Biểu đồ use case phân rã chức năng quản lý khen thưởng}
    \label{fig:usecase-dexuat}
\end{figure}

\subsection{Chức năng xét danh hiệu hằng năm cá nhân}
\label{subsection:2.2.5}

Đây là nhóm chức năng có nhiều quy tắc xét điều kiện đan xen.
Quy tắc danh hiệu hằng năm đi kèm điều kiện tính theo khoảng năm gần nhất tính lùi, lại có một danh hiệu chỉ được nhận đúng một lần trong toàn bộ thời gian phục vụ.
Các danh hiệu cao hơn của cá nhân xét trên chuỗi CSTĐCS liên tục gồm ba cấp BKBQP, CSTĐTQ và BKTTCP với các điều kiện sau:

\begin{itemize}
    \item \textbf{BKBQP}: chu kỳ 2 năm, yêu cầu 2 năm CSTĐCS liên tục và NCKH mỗi năm; không có danh hiệu tiền điều kiện; có thể nhận lặp lại sau mỗi chu kỳ.
    \item \textbf{CSTĐTQ}: chu kỳ 3 năm, yêu cầu 3 năm CSTĐCS liên tục, có ít nhất một BKBQP trong ba năm gần nhất tính từ năm liền trước năm xét, kèm NCKH mỗi năm, và có thể nhận lặp lại sau mỗi chu kỳ. Vì khoảng thời gian này dịch dần theo năm xét, BKBQP của chu kỳ trước không còn nằm trong khoảng đang xét, nên người được xét phải có một BKBQP mới thuộc khoảng đó.
    \item \textbf{BKTTCP}: chu kỳ 7 năm, yêu cầu 7 năm CSTĐCS liên tục với độ dài chuỗi là bội số của 7, đúng 3 BKBQP và đúng 2 CSTĐTQ nằm trong khoảng bảy năm gần nhất, kèm NCKH mỗi năm. Với cá nhân, danh hiệu này chỉ được xét cấp một lần, không lặp lại theo chu kỳ.
\end{itemize}

Use case chính của nhóm chức năng này là \textbf{Xem gợi ý khen thưởng}, đã xuất hiện trong biểu đồ quản lý khen thưởng ở Hình~\ref{fig:usecase-dexuat}.
Hệ thống tính từ dữ liệu lịch sử rồi sinh ra hai loại thông tin cho mỗi quân nhân.
Thứ nhất là danh sách danh hiệu đủ điều kiện đề nghị trong năm hiện tại kèm lý do.
Thứ hai là các chỉ số theo khoảng thời gian xét, gồm số BKBQP trong khoảng 3 năm hoặc 7 năm gần nhất, số CSTĐTQ trong khoảng 7 năm gần nhất và độ dài chuỗi CSTĐCS liên tục đang đếm.
Cán bộ Phòng Chính trị tham chiếu các gợi ý này để chốt danh sách đề xuất.

Đi kèm chức năng này là cơ chế tính lại điều kiện tự động: mỗi khi một bản ghi danh hiệu hằng năm thay đổi, hệ thống tính lại điều kiện cho riêng quân nhân bị ảnh hưởng; việc tính lại cho toàn bộ quân nhân được chạy tự động theo lịch định kỳ, và Cán bộ Phòng Chính trị cũng có thể kích hoạt thủ công khi cần. Cơ chế này giữ cho hồ sơ điều kiện hằng năm luôn nhất quán với dữ liệu danh hiệu đã ghi nhận.

\subsection{Chức năng xét danh hiệu hằng năm đơn vị}
\label{subsection:2.2.6}

Danh hiệu hằng năm đơn vị có cấu trúc đơn giản hơn cá nhân do không có cấp CSTĐTQ và không yêu cầu thành tích NCKH.
Các danh hiệu cao hơn của đơn vị xét trên chuỗi ĐVQT liên tục chỉ gồm hai cấp:

\begin{itemize}
    \item \textbf{BKBQP đơn vị} chu kỳ 2 năm, yêu cầu 2 năm ĐVQT liên tục, không có danh hiệu tiền điều kiện, có thể nhận lặp lại sau mỗi chu kỳ.
    \item \textbf{BKTTCP đơn vị} chu kỳ 7 năm, yêu cầu 7 năm ĐVQT liên tục và ít nhất 3 BKBQP đơn vị trong 7 năm gần nhất, có thể nhận lặp lại sau mỗi chu kỳ.
\end{itemize}

Điểm khác với cá nhân là đơn vị có thể được xét lại BKTTCP sau mỗi chu kỳ, nên yêu cầu về các danh hiệu tiền điều kiện cũng ít chặt chẽ hơn so với cá nhân.

\subsection{Biểu đồ use case phân rã chức năng nhập và xuất dữ liệu Excel}
\label{subsection:2.2.7}

Lượng dữ liệu lịch sử của Học viện hiện lưu trên các tài liệu và bản giấy.
Chức năng nhập Excel hàng loạt cung cấp tệp mẫu chuẩn để cán bộ chuyển dữ liệu cũ vào hệ thống thay vì nhập tay từng bản ghi.
Quy trình gồm bốn use case con, thể hiện ở Hình~\ref{fig:uc-excel}:

\begin{itemize}
    \item \textbf{Tải tệp mẫu}: hệ thống sinh tệp Excel có cấu trúc cột chuẩn theo từng loại khen thưởng (danh hiệu hằng năm, khen thưởng đơn vị hằng năm, HCCSVV, HCBVTQ, ...) cùng các ràng buộc kiểm tra dữ liệu cho phép chọn từ danh sách thả xuống ngay trong tệp.
    \item \textbf{Xem trước}: sau khi cán bộ điền dữ liệu và tải tệp lên, hệ thống đọc nội dung từng dòng, đối chiếu với quy tắc kiểm tra dữ liệu và các quy tắc nghiệp vụ, sau đó hiển thị kết quả gồm các dòng hợp lệ và các dòng có lỗi kèm chỉ số dòng và mô tả lỗi. Bước này không ghi vào cơ sở dữ liệu.
    \item \textbf{Xác nhận}: cán bộ kiểm tra danh sách hợp lệ rồi xác nhận. Hệ thống ghi tuần tự từng dòng; nếu xảy ra lỗi tại bất kỳ dòng nào, toàn bộ thao tác được hủy bỏ và cơ sở dữ liệu giữ nguyên trạng thái trước khi nhập.
    \item \textbf{Xuất danh sách}: cho phép xuất kết quả khen thưởng theo nhiều tiêu chí (theo năm, theo đơn vị, theo loại danh hiệu) ra tệp Excel.
\end{itemize}

\begin{figure}[H]
    \centering
    \includegraphics[width=\textwidth,height=0.9\textheight,keepaspectratio]{uc-excel}
    \caption{Biểu đồ use case phân rã chức năng nhập và xuất dữ liệu Excel}
    \label{fig:uc-excel}
\end{figure}

\subsection{Biểu đồ use case phân rã chức năng quản lý quyết định khen thưởng}
\label{subsection:2.2.8a}

Mọi danh hiệu, huân/huy chương sau khi được phê duyệt đều gắn với một quyết định khen thưởng.
Một quyết định lưu số quyết định duy nhất, ngày ký, người ký, tệp PDF đính kèm và loại khen thưởng.
Quyết định này được tham chiếu bởi tám bảng khen thưởng qua số quyết định.
Nhóm chức năng này do Cán bộ Phòng Chính trị quản lý, gồm các use case con: xem danh sách quyết định, tìm kiếm và lọc, tải lên tệp PDF quyết định, cập nhật thông tin, tải xuống tệp PDF và xoá.
Thao tác xoá có thêm bước kiểm tra ràng buộc.
Nếu quyết định đang được ít nhất một bản ghi khen thưởng tham chiếu thì hệ thống chặn xoá để bảo toàn tính toàn vẹn liên kết giữa quyết định và các danh hiệu đã trao.
Các use case của nhóm chức năng quản lý quyết định được thể hiện ở Hình~\ref{fig:usecase-quyetdinh}.

\begin{figure}[H]
    \centering
    \includegraphics[width=\textwidth,height=0.9\textheight,keepaspectratio]{uc-quyet-dinh}
    \caption{Biểu đồ use case phân rã chức năng quản lý quyết định khen thưởng}
    \label{fig:usecase-quyetdinh}
\end{figure}

\subsection{Biểu đồ use case phân rã chức năng quản lý tài khoản}
\label{subsection:2.2.8}

Nhóm chức năng quản lý tài khoản phục vụ tạo lập và duy trì tài khoản người dùng, gồm:

\begin{itemize}
    \item \textbf{Quản lý tài khoản}: Quản trị viên tạo được tài khoản ở mọi vai trò, còn Cán bộ Phòng Chính trị chỉ tạo tài khoản Chỉ huy đơn vị và Người dùng; mỗi tài khoản gắn với một quân nhân duy nhất.
    \item \textbf{Đặt lại mật khẩu}: Cán bộ Phòng Chính trị có thể đặt lại mật khẩu cho Chỉ huy đơn vị / Người dùng về giá trị mặc định.
\end{itemize}

Tất cả các thao tác trong nhóm chức năng này đều có lưu nhật ký riêng để đảm bảo trách nhiệm cá nhân và khả năng truy vết.
Các use case quản lý tài khoản được thể hiện ở Hình~\ref{fig:usecase-quantri}. Các chức năng quản trị còn lại như nhật ký hệ thống và sao lưu đã xuất hiện trong biểu đồ tổng quát ở Hình~\ref{fig:usecase-tongquat} và được đặc tả chi tiết ở Mục~\ref{section:2.3}.

\begin{figure}[H]
    \centering
    \includegraphics[width=\textwidth,height=0.9\textheight,keepaspectratio]{uc-tai-khoan}
    \caption{Biểu đồ use case phân rã chức năng quản lý tài khoản}
    \label{fig:usecase-quantri}
\end{figure}

\subsection{Ba quy trình nghiệp vụ quan trọng}
\label{subsection:2.2.9}

Bên cạnh biểu đồ use case, ba quy trình nghiệp vụ tiêu biểu được mô hình hoá thêm bằng biểu đồ hoạt động (activity diagram) để làm rõ luồng xử lý từng bước.

\textbf{Quy trình A: Đề xuất và phê duyệt một đợt khen thưởng.} Toàn bộ luồng đề xuất và phê duyệt được thể hiện ở Hình~\ref{fig:activity-quytrinh}.

\begin{figure}[H]
    \centering
    \includegraphics[width=\textwidth,height=0.9\textheight,keepaspectratio]{activity-quy-trinh}
    \caption{Biểu đồ hoạt động quy trình đề xuất và phê duyệt khen thưởng}
    \label{fig:activity-quytrinh}
\end{figure}

Chỉ huy đơn vị chọn loại đề xuất và năm xét, hệ thống lọc danh sách quân nhân/đơn vị thuộc phạm vi quản lý.
Chỉ huy đơn vị chọn các đối tượng đủ điều kiện, thiết lập danh hiệu cụ thể rồi gửi đề xuất.
Hệ thống lưu đề xuất ở trạng thái chờ duyệt và gửi thông báo thời gian thực tới Cán bộ Phòng Chính trị.
Cán bộ Phòng Chính trị mở danh sách chờ duyệt, đối chiếu từng đối tượng và bổ sung số quyết định trước khi phê duyệt; nếu đề xuất chưa phù hợp, cán bộ từ chối kèm lý do để Chỉ huy đơn vị điều chỉnh và đề xuất lại.
Khi phê duyệt, hệ thống thực hiện đồng thời một loạt thao tác: đổi trạng thái đề xuất thành đã duyệt, lưu các bản ghi khen thưởng vào đúng bảng theo loại danh hiệu kèm số quyết định và tệp PDF, ghi nhật ký, thông báo cho Chỉ huy đơn vị đã đề xuất và tính lại các hồ sơ điều kiện liên quan.
Cả loạt thao tác này hoặc cùng hoàn tất, hoặc cùng bị huỷ để dữ liệu luôn nhất quán.

\textbf{Quy trình B: Nhập tệp Excel hàng loạt.} Luồng nhập Excel hai bước được thể hiện ở Hình~\ref{fig:activity-excel}.

\begin{figure}[H]
    \centering
    \includegraphics[width=\textwidth,height=0.9\textheight,keepaspectratio]{activity-excel}
    \caption{Biểu đồ hoạt động quy trình nhập khen thưởng từ tệp Excel}
    \label{fig:activity-excel}
\end{figure}

Cán bộ Phòng Chính trị tải lên tệp Excel chứa danh sách bản ghi cần nhập.
Hệ thống đọc tệp theo loại khen thưởng và xử lý từng dòng.
Mỗi dòng được kiểm tra theo các quy tắc nghiệp vụ liên quan, chẳng hạn tìm được quân nhân theo số căn cước và đơn vị phải tồn tại.
Các dòng hợp lệ và lỗi được tổng hợp thành một kết quả xem trước hiển thị cho cán bộ.
Cán bộ xem trước, đối chiếu, có thể quay lại sửa tệp Excel rồi tải lên lại.
Khi xác nhận, hệ thống ghi tuần tự các dòng hợp lệ; nếu có bất kỳ lỗi nào tại thời điểm ghi, toàn bộ thao tác bị huỷ bỏ và cơ sở dữ liệu giữ nguyên như trước.

\textbf{Quy trình C: Thêm khen thưởng đột xuất.} Luồng thêm khen thưởng đột xuất được thể hiện ở Hình~\ref{fig:activity-dotxuat}.

\begin{figure}[H]
    \centering
    \includegraphics[width=\textwidth,height=0.9\textheight,keepaspectratio]{activity-dot-xuat}
    \caption{Biểu đồ hoạt động quy trình thêm khen thưởng đột xuất}
    \label{fig:activity-dotxuat}
\end{figure}

Khen thưởng đột xuất phát sinh ngoài chu kỳ và đã có quyết định của cấp trên, nên Cán bộ Phòng Chính trị thêm trực tiếp vào hệ thống mà không qua bước đề xuất và phê duyệt.
Cán bộ chọn đối tượng là cá nhân hay tập thể, nhập năm và hình thức khen thưởng, chọn danh sách quân nhân hoặc đơn vị được khen, đính kèm tệp nếu có rồi gắn số quyết định, có thể chọn lại một quyết định đã có hoặc tạo quyết định mới nếu chưa có trong hệ thống.
Sau khi cán bộ xác nhận, hệ thống kiểm tra dữ liệu đầu vào, lưu bản ghi khen thưởng, gửi thông báo tới đơn vị và quân nhân liên quan rồi ghi nhật ký hệ thống.

\section{Đặc tả use case chi tiết}
\label{section:2.3}

\subsection{Đặc tả use case đăng nhập}

\begin{table}[H]
\centering
\renewcommand{\arraystretch}{1.1}
\begin{tabular}{|>{\centering\arraybackslash}m{3.2cm}|>{\centering\arraybackslash}m{0.8cm}|>{\centering\arraybackslash}m{2.8cm}|>{\raggedright\arraybackslash}m{6.2cm}|}
\hline
\textbf{Tên use case} & \multicolumn{3}{m{10.4cm}|}{Đăng nhập hệ thống} \\ \hline
\textbf{Tác nhân} & \multicolumn{3}{m{10.4cm}|}{Quản trị viên, Cán bộ Phòng Chính trị, Chỉ huy đơn vị, Người dùng} \\ \hline
\textbf{Tiền điều kiện} & \multicolumn{3}{m{10.4cm}|}{Người dùng đã có tài khoản hợp lệ.} \\ \hline
\multirow{7}{*}{\parbox{3.2cm}{\centering\textbf{Luồng sự kiện chính}\\\textbf{(Thành công)}}}
 & \textbf{STT} & \textbf{Thực hiện bởi} & \multicolumn{1}{c|}{\textbf{Hành động}} \\ \cline{2-4}
 & 1. & Người dùng & Truy cập trang đăng nhập \\ \cline{2-4}
 & 2. & Người dùng & Nhập tên đăng nhập và mật khẩu \\ \cline{2-4}
 & 3. & Người dùng & Bấm nút ``Đăng nhập'' \\ \cline{2-4}
 & 4. & Hệ thống & Kiểm tra tên đăng nhập và mật khẩu \\ \cline{2-4}
 & 5. & Hệ thống & Tạo phiên đăng nhập cho người dùng \\ \cline{2-4}
 & 6. & Hệ thống & Chuyển hướng tới trang chủ tương ứng vai trò \\ \hline
\multirow{3}{*}{\parbox{3cm}{\centering\textbf{Luồng sự kiện}\\\textbf{thay thế}}}
 & \textbf{STT} & \textbf{Thực hiện bởi} & \multicolumn{1}{c|}{\textbf{Hành động}} \\ \cline{2-4}
 & 4a. & Hệ thống & Báo ``Tên đăng nhập hoặc mật khẩu không đúng'' \\ \cline{2-4}
 & 4b. & Hệ thống & Quá số lần đăng nhập sai cho phép, hệ thống tạm thời từ chối thêm yêu cầu từ địa chỉ đó trong ít phút \\ \hline
\textbf{Hậu điều kiện} & \multicolumn{3}{m{10.4cm}|}{Phiên đăng nhập được thiết lập; người dùng vào trang chủ tương ứng vai trò.} \\ \hline
\end{tabular}
\caption{Đặc tả use case đăng nhập}
\label{tab:uc01}
\end{table}

\begin{table}[H]
\centering
\renewcommand{\arraystretch}{1.1}
\begin{tabular}{|>{\centering\arraybackslash}m{0.8cm}|>{\centering\arraybackslash}m{3.0cm}|>{\raggedright\arraybackslash}m{4.2cm}|>{\centering\arraybackslash}m{1.8cm}|>{\centering\arraybackslash}m{2.8cm}|}
\hline
\textbf{STT} & \textbf{Trường dữ liệu} & \multicolumn{1}{c|}{\textbf{Mô tả}} & \textbf{Bắt buộc} & \textbf{Ví dụ hợp lệ} \\ \hline
1 & Tên đăng nhập & Tên tài khoản đăng nhập & Có & \texttt{ducta} \\ \hline
2 & Mật khẩu & Mật khẩu đăng nhập & Có & \texttt{Dat\#2026} \\ \hline
\end{tabular}
\caption{Dữ liệu đầu vào cho chức năng đăng nhập}
\label{tab:uc01-input}
\end{table}

\subsection{Đặc tả use case quản lý lịch sử chức vụ quân nhân}

\begin{table}[H]
\centering
\renewcommand{\arraystretch}{1.1}
\begin{tabular}{|>{\centering\arraybackslash}m{3.2cm}|>{\centering\arraybackslash}m{0.8cm}|>{\centering\arraybackslash}m{2.8cm}|>{\raggedright\arraybackslash}m{6.2cm}|}
\hline
\textbf{Tên use case} & \multicolumn{3}{m{10.4cm}|}{Quản lý lịch sử chức vụ quân nhân} \\ \hline
\textbf{Tác nhân} & \multicolumn{3}{m{10.4cm}|}{Cán bộ Phòng Chính trị, Chỉ huy đơn vị} \\ \hline
\textbf{Tiền điều kiện} & \multicolumn{3}{m{10.4cm}|}{Đã đăng nhập (Cán bộ Phòng Chính trị hoặc Chỉ huy đơn vị); đã có dữ liệu quân nhân và chức vụ. Chỉ huy đơn vị chỉ thao tác trong đơn vị mình.} \\ \hline
\multirow{7}{*}{\parbox{3.2cm}{\centering\textbf{Luồng sự kiện chính}\\\textbf{(Thành công)}}}
 & \textbf{STT} & \textbf{Thực hiện bởi} & \multicolumn{1}{c|}{\textbf{Hành động}} \\ \cline{2-4}
 & 1. & Cán bộ Phòng Chính trị & Mở hồ sơ quân nhân, vào mục lịch sử chức vụ và bấm ``Thêm lịch sử'' \\ \cline{2-4}
 & 2. & Cán bộ Phòng Chính trị & Chọn chức vụ và nhập ngày bắt đầu; nhập thêm ngày kết thúc nếu là giai đoạn đã qua \\ \cline{2-4}
 & 3. & Cán bộ Phòng Chính trị & Bấm ``Lưu'' \\ \cline{2-4}
 & 4. & Hệ thống & Kiểm tra chức vụ tồn tại, ngày bắt đầu và ngày kết thúc hợp lệ \\ \cline{2-4}
 & 5. & Hệ thống & Kiểm tra không trùng thời gian với các giai đoạn đã có; tính số tháng đảm nhiệm rồi lưu kèm hệ số chức vụ \\ \cline{2-4}
 & 6. & Hệ thống & Ghi nhật ký hệ thống và trả kết quả thành công \\ \hline
\multirow{5}{*}{\parbox{3cm}{\centering\textbf{Luồng sự kiện}\\\textbf{thay thế}}}
 & \textbf{STT} & \textbf{Thực hiện bởi} & \multicolumn{1}{c|}{\textbf{Hành động}} \\ \cline{2-4}
 & 1a. & Cán bộ Phòng Chính trị & Thay vì thêm mới, chọn một giai đoạn đã có để sửa thông tin hoặc xoá khỏi danh sách; hệ thống áp dụng lại các kiểm tra ở bước 4--5 \\ \cline{2-4}
 & 2a. & Chỉ huy đơn vị & Nhập giai đoạn trùng thời gian với giai đoạn hiện tại: hệ thống báo lỗi; giai đoạn hiện tại không đặt được ngày kết thúc (chỉ Phòng Chính trị đổi chức vụ hiện tại) \\ \cline{2-4}
 & 4a. & Hệ thống & Ngày kết thúc trước ngày bắt đầu hoặc thiếu trường bắt buộc: báo lỗi và yêu cầu nhập lại \\ \cline{2-4}
 & 5a. & Hệ thống & Giai đoạn nhập vào trùng thời gian với một giai đoạn chức vụ đã có: hệ thống chặn và yêu cầu nhập lại \\ \hline
\textbf{Hậu điều kiện} & \multicolumn{3}{m{10.4cm}|}{Bảng lịch sử chức vụ được cập nhật; mỗi giai đoạn lưu hệ số và số tháng đảm nhiệm — đầu vào để xét HCBVTQ.} \\ \hline
\end{tabular}
\caption{Đặc tả use case quản lý lịch sử chức vụ quân nhân}
\label{tab:uc02}
\end{table}

\begin{table}[H]
\centering
\renewcommand{\arraystretch}{1.1}
\begin{tabular}{|>{\centering\arraybackslash}m{0.8cm}|>{\centering\arraybackslash}m{3.0cm}|>{\raggedright\arraybackslash}m{4.2cm}|>{\centering\arraybackslash}m{1.8cm}|>{\centering\arraybackslash}m{2.8cm}|}
\hline
\textbf{STT} & \textbf{Trường dữ liệu} & \multicolumn{1}{c|}{\textbf{Mô tả}} & \textbf{Bắt buộc} & \textbf{Ví dụ hợp lệ} \\ \hline
1 & Chức vụ & Chức vụ quân nhân giữ trong giai đoạn & Có & Trợ lý \\ \hline
2 & Ngày bắt đầu & Ngày bắt đầu giữ chức vụ & Có & 01/01/2020 \\ \hline
3 & Ngày kết thúc & Để trống nếu là giai đoạn hiện tại; bắt buộc với giai đoạn đã qua & Khi đã qua & 31/12/2023 \\ \hline
\end{tabular}
\caption{Dữ liệu đầu vào cho chức năng quản lý lịch sử chức vụ}
\label{tab:uc02-input}
\end{table}

\subsection{Đặc tả use case tạo đề xuất khen thưởng}

\begin{table}[H]
\centering
\renewcommand{\arraystretch}{1.1}
\begin{tabular}{|>{\centering\arraybackslash}m{3.2cm}|>{\centering\arraybackslash}m{0.8cm}|>{\centering\arraybackslash}m{2.8cm}|>{\raggedright\arraybackslash}m{6.2cm}|}
\hline
\textbf{Tên use case} & \multicolumn{3}{m{10.4cm}|}{Tạo đề xuất khen thưởng cho một đợt} \\ \hline
\textbf{Tác nhân} & \multicolumn{3}{m{10.4cm}|}{Chỉ huy đơn vị} \\ \hline
\textbf{Tiền điều kiện} & \multicolumn{3}{m{10.4cm}|}{Đã đăng nhập; có ít nhất một quân nhân thuộc phạm vi quản lý của Chỉ huy đơn vị.} \\ \hline
\multirow{11}{*}{\parbox{3.2cm}{\centering\textbf{Luồng sự kiện chính}\\\textbf{(Thành công)}}}
 & \textbf{STT} & \textbf{Thực hiện bởi} & \multicolumn{1}{c|}{\textbf{Hành động}} \\ \cline{2-4}
 & 1. & Chỉ huy đơn vị & Truy cập ``Đề xuất khen thưởng'' và bấm ``Tạo đề xuất mới'' \\ \cline{2-4}
 & 2. & Chỉ huy đơn vị & Chọn loại đề xuất \\ \cline{2-4}
 & 3. & Chỉ huy đơn vị & Chọn năm/tháng xét của đợt khen thưởng \\ \cline{2-4}
 & 4. & Hệ thống & Lọc và hiển thị danh sách quân nhân/đơn vị đủ điều kiện \\ \cline{2-4}
 & 5. & Chỉ huy đơn vị & Tích chọn các đối tượng đề xuất \\ \cline{2-4}
 & 6. & Chỉ huy đơn vị & Thiết lập danh hiệu cụ thể cho từng đối tượng \\ \cline{2-4}
 & 7. & Chỉ huy đơn vị & Tải lên các tệp đính kèm (nếu có) \\ \cline{2-4}
 & 8. & Chỉ huy đơn vị & Xem lại toàn bộ thông tin và bấm ``Gửi đề xuất'' \\ \cline{2-4}
 & 9. & Hệ thống & Xác thực dữ liệu, tạo bản ghi đề xuất ở trạng thái chờ duyệt \\ \cline{2-4}
 & 10. & Hệ thống & Gửi thông báo thời gian thực tới Cán bộ Phòng Chính trị \\ \hline
\multirow{3}{*}{\parbox{3cm}{\centering\textbf{Luồng sự kiện}\\\textbf{thay thế}}}
 & \textbf{STT} & \textbf{Thực hiện bởi} & \multicolumn{1}{c|}{\textbf{Hành động}} \\ \cline{2-4}
 & 5a. & Hệ thống & Có đối tượng chưa đủ điều kiện, hệ thống từ chối và báo lỗi, không cho thêm \\ \cline{2-4}
 & 8a. & Hệ thống & Bản nháp được lưu tạm trên trình duyệt; khi rời trang hoặc gửi lỗi, dữ liệu không bị mất \\ \hline
\textbf{Hậu điều kiện} & \multicolumn{3}{m{10.4cm}|}{Bản ghi đề xuất được tạo ở trạng thái chờ duyệt; thông báo đã gửi tới Cán bộ Phòng Chính trị; nhật ký hệ thống ghi lại thao tác tạo đề xuất cùng tóm tắt loại và số lượng đối tượng.} \\ \hline
\end{tabular}
\caption{Đặc tả use case tạo đề xuất khen thưởng}
\label{tab:uc03}
\end{table}

\begin{table}[H]
\centering
\renewcommand{\arraystretch}{1.1}
\begin{tabular}{|>{\centering\arraybackslash}m{0.8cm}|>{\centering\arraybackslash}m{3.0cm}|>{\raggedright\arraybackslash}m{4.2cm}|>{\centering\arraybackslash}m{1.8cm}|>{\centering\arraybackslash}m{2.8cm}|}
\hline
\textbf{STT} & \textbf{Trường dữ liệu} & \multicolumn{1}{c|}{\textbf{Mô tả}} & \textbf{Bắt buộc} & \textbf{Ví dụ hợp lệ} \\ \hline
1 & Loại đề xuất & Một trong bảy nhóm khen thưởng & Có & HCCSVV \\ \hline
2 & Năm xét & Năm xét khen thưởng & Có & 2026 \\ \hline
3 & Tháng nhận & Chỉ nhập với loại xét theo mốc thời gian phục vụ (HCCSVV, HCBVTQ, HCQKQT, KNC); loại hằng năm chỉ cần năm & Tuỳ loại & 6 \\ \hline
4 & Danh sách đối tượng & Quân nhân hoặc đơn vị được đề xuất & Có & 25 quân nhân \\ \hline
5 & Danh hiệu cụ thể & Hạng/cấp danh hiệu cho từng đối tượng & Có & Hạng Ba \\ \hline
6 & Tệp đính kèm & Tệp PDF biên bản, ảnh giấy khen & Không & \texttt{bb-\allowbreak{}2026.pdf} \\ \hline
\end{tabular}
\caption{Dữ liệu đầu vào cho chức năng tạo đề xuất khen thưởng}
\label{tab:uc03-input}
\end{table}

\subsection{Đặc tả use case phê duyệt đề xuất}

\begin{table}[H]
\centering
\renewcommand{\arraystretch}{1.1}
\begin{tabular}{|>{\centering\arraybackslash}m{3.2cm}|>{\centering\arraybackslash}m{0.8cm}|>{\centering\arraybackslash}m{2.8cm}|>{\raggedright\arraybackslash}m{6.2cm}|}
\hline
\textbf{Tên use case} & \multicolumn{3}{m{10.4cm}|}{Phê duyệt một đề xuất khen thưởng} \\ \hline
\textbf{Tác nhân} & \multicolumn{3}{m{10.4cm}|}{Cán bộ Phòng Chính trị} \\ \hline
\textbf{Tiền điều kiện} & \multicolumn{3}{m{10.4cm}|}{Đề xuất đang ở trạng thái chờ duyệt; Cán bộ Phòng Chính trị có quyền phê duyệt loại đề xuất tương ứng.} \\ \hline
\multirow{8}{*}{\parbox{3.2cm}{\centering\textbf{Luồng sự kiện chính}\\\textbf{(Thành công)}}}
 & \textbf{STT} & \textbf{Thực hiện bởi} & \multicolumn{1}{c|}{\textbf{Hành động}} \\ \cline{2-4}
 & 1. & Cán bộ Phòng Chính trị & Mở chi tiết đề xuất từ danh sách ``Chờ duyệt'' \\ \cline{2-4}
 & 2. & Cán bộ Phòng Chính trị & Đối chiếu danh sách đối tượng và danh hiệu \\ \cline{2-4}
 & 3. & Cán bộ Phòng Chính trị & Nhập số quyết định đã có hoặc thêm mới số quyết định; với loại xét theo mốc thời gian phục vụ thì nhập thêm tháng, năm nhận \\ \cline{2-4}
 & 4. & Cán bộ Phòng Chính trị & Bấm ``Phê duyệt'' \\ \cline{2-4}
 & 5. & Hệ thống & Ghi danh hiệu, chuyển trạng thái đã duyệt và lưu số quyết định; sau khi lưu xong, tính lại hồ sơ điều kiện \\ \cline{2-4}
 & 6. & Hệ thống & Gửi thông báo thời gian thực tới Chỉ huy đơn vị đã đề xuất và các quân nhân được khen thưởng \\ \cline{2-4}
 & 7. & Hệ thống & Ghi nhật ký hành động phê duyệt kèm mô tả chi tiết \\ \hline
\multirow{5}{*}{\parbox{3cm}{\centering\textbf{Luồng sự kiện}\\\textbf{thay thế}}}
 & \textbf{STT} & \textbf{Thực hiện bởi} & \multicolumn{1}{c|}{\textbf{Hành động}} \\ \cline{2-4}
 & 3a. & Cán bộ Phòng Chính trị & Phát hiện danh hiệu đề xuất chưa phù hợp, chuyển sang luồng từ chối kèm lý do (5a); Chỉ huy đơn vị điều chỉnh và đề xuất lại \\ \cline{2-4}
 & 4a. & Hệ thống & Thiếu số quyết định, hệ thống chặn và yêu cầu nhập đủ \\ \cline{2-4}
 & 5a. & Cán bộ Phòng Chính trị & Từ chối thay vì phê duyệt: trạng thái chuyển ``Từ chối'' kèm lý do \\ \cline{2-4}
 & 5b. & Hệ thống & Gặp lỗi hệ thống khi ghi: huỷ toàn bộ giao dịch, giữ nguyên dữ liệu \\ \hline
\textbf{Hậu điều kiện} & \multicolumn{3}{m{10.4cm}|}{Trạng thái đề xuất là ``Đã duyệt'' hoặc ``Từ chối''; nếu được duyệt, các bản ghi khen thưởng đã có trong cơ sở dữ liệu, hồ sơ điều kiện được tính lại, thông báo đã gửi, nhật ký đầy đủ.} \\ \hline
\end{tabular}
\caption{Đặc tả use case phê duyệt đề xuất}
\label{tab:uc04}
\end{table}

\begin{table}[H]
\centering
\renewcommand{\arraystretch}{1.1}
\begin{tabular}{|>{\centering\arraybackslash}m{0.8cm}|>{\centering\arraybackslash}m{3.0cm}|>{\raggedright\arraybackslash}m{4.2cm}|>{\centering\arraybackslash}m{1.8cm}|>{\centering\arraybackslash}m{2.8cm}|}
\hline
\textbf{STT} & \textbf{Trường dữ liệu} & \multicolumn{1}{c|}{\textbf{Mô tả}} & \textbf{Bắt buộc} & \textbf{Ví dụ hợp lệ} \\ \hline
1 & Số quyết định & Số quyết định khen thưởng (tham chiếu PDF đã tải lên) & Có & 1234/QĐ-HV \\ \hline
2 & Tháng và năm nhận & Bắt buộc với loại theo mốc thời gian (HCCSVV, HCBVTQ, HCQKQT, KNC); loại hằng năm và NCKH chỉ cần năm & Tùy loại & 12/2025 \\ \hline
3 & Lý do từ chối & Áp dụng khi từ chối đề xuất & Khi từ chối & Hạng đề xuất chưa phù hợp \\ \hline
\end{tabular}
\caption{Dữ liệu đầu vào cho chức năng phê duyệt đề xuất}
\label{tab:uc04-input}
\end{table}

\subsection{Đặc tả use case tính lại điều kiện khen thưởng và tạo gợi ý}

\begin{table}[H]
\centering
\renewcommand{\arraystretch}{1.1}
\begin{tabular}{|>{\centering\arraybackslash}m{3.2cm}|>{\centering\arraybackslash}m{0.8cm}|>{\centering\arraybackslash}m{2.8cm}|>{\raggedright\arraybackslash}m{6.2cm}|}
\hline
\textbf{Tên use case} & \multicolumn{3}{m{10.4cm}|}{Tính lại điều kiện khen thưởng hằng năm và tạo gợi ý} \\ \hline
\textbf{Tác nhân} & \multicolumn{3}{m{10.4cm}|}{Cán bộ Phòng Chính trị, Chỉ huy đơn vị (kích hoạt khi thêm hoặc đề xuất khen thưởng); Hệ thống (tự động sau mỗi lần phê duyệt).} \\ \hline
\textbf{Tiền điều kiện} & \multicolumn{3}{m{10.4cm}|}{Quân nhân tồn tại; có ít nhất một bản ghi danh hiệu hằng năm.} \\ \hline
\multirow{6}{*}{\parbox{3.2cm}{\centering\textbf{Luồng sự kiện chính}\\\textbf{(Thành công)}}}
 & \textbf{STT} & \textbf{Thực hiện bởi} & \multicolumn{1}{c|}{\textbf{Hành động}} \\ \cline{2-4}
 & 1. & Hệ thống & Truy vấn toàn bộ danh hiệu hằng năm của quân nhân \\ \cline{2-4}
 & 2. & Hệ thống & Tính ngữ cảnh danh hiệu hằng năm (số năm liên tục, danh hiệu trong khoảng thời gian xét) \\ \cline{2-4}
 & 3. & Hệ thống & Với mỗi cấp danh hiệu, gọi hàm xét điều kiện \\ \cline{2-4}
 & 4. & Hệ thống & Ghi kết quả vào bảng hồ sơ điều kiện \\ \cline{2-4}
 & 5. & Hệ thống & Trả nội dung gợi ý cho giao diện hiển thị \\ \hline
\multirow{4}{*}{\parbox{3cm}{\centering\textbf{Luồng sự kiện}\\\textbf{thay thế}}}
 & \textbf{STT} & \textbf{Thực hiện bởi} & \multicolumn{1}{c|}{\textbf{Hành động}} \\ \cline{2-4}
 & 3a. & Hệ thống & Quân nhân đã nhận BKTTCP nên hệ thống không gợi ý thêm, vì chưa hỗ trợ danh hiệu cao hơn \\ \cline{2-4}
 & 3b. & Hệ thống & Số năm CSTĐCS liên tục chưa đủ chu kỳ, gợi ý chưa đủ điều kiện xét danh hiệu \\ \cline{2-4}
 & 3c. & Hệ thống & Quân nhân lỡ một hoặc nhiều đợt đề nghị; chu kỳ vẫn đếm tiếp và hệ thống đề nghị lại ở chu kỳ kế tiếp \\ \hline
\textbf{Hậu điều kiện} & \multicolumn{3}{m{10.4cm}|}{Hồ sơ điều kiện nhất quán với dữ liệu nguồn; trang chi tiết quân nhân hiển thị các danh hiệu đủ điều kiện kèm nội dung gợi ý.} \\ \hline
\end{tabular}
\caption{Đặc tả use case tính lại điều kiện khen thưởng và tạo gợi ý}
\label{tab:uc05}
\end{table}

\subsection{Đặc tả use case nhập dữ liệu khen thưởng hàng loạt từ tệp Excel}

\begin{table}[H]
\centering
\renewcommand{\arraystretch}{1.1}
\begin{tabular}{|>{\centering\arraybackslash}m{3.2cm}|>{\centering\arraybackslash}m{0.8cm}|>{\centering\arraybackslash}m{2.8cm}|>{\raggedright\arraybackslash}m{6.2cm}|}
\hline
\textbf{Tên use case} & \multicolumn{3}{m{10.4cm}|}{Nhập dữ liệu khen thưởng hàng loạt từ tệp Excel} \\ \hline
\textbf{Tác nhân} & \multicolumn{3}{m{10.4cm}|}{Cán bộ Phòng Chính trị} \\ \hline
\textbf{Tiền điều kiện} & \multicolumn{3}{m{10.4cm}|}{Đã có tệp mẫu tải xuống từ hệ thống và điền dữ liệu theo cấu trúc cột quy định.} \\ \hline
\multirow{9}{*}{\parbox{3.2cm}{\centering\textbf{Luồng sự kiện chính}\\\textbf{(Thành công)}}}
 & \textbf{STT} & \textbf{Thực hiện bởi} & \multicolumn{1}{c|}{\textbf{Hành động}} \\ \cline{2-4}
 & 1. & Cán bộ Phòng Chính trị & Truy cập chức năng thêm khen thưởng và chọn loại khen thưởng \\ \cline{2-4}
 & 2. & Cán bộ Phòng Chính trị & Tải lên tệp Excel qua biểu mẫu \\ \cline{2-4}
 & 3. & Hệ thống & Đọc tệp, kiểm tra cấu trúc cột và đọc tuần tự từng dòng \\ \cline{2-4}
 & 4. & Hệ thống & Kiểm tra và tìm quân nhân tương ứng cho mỗi dòng \\ \cline{2-4}
 & 5. & Hệ thống & Trả về bảng xem trước với các dòng hợp lệ và dòng lỗi kèm mô tả \\ \cline{2-4}
 & 6. & Cán bộ Phòng Chính trị & Kiểm tra danh sách lỗi, nếu đồng ý bấm ``Xác nhận nhập'' \\ \cline{2-4}
 & 7. & Hệ thống & Ghi tuần tự các dòng hợp lệ vào cơ sở dữ liệu \\ \cline{2-4}
 & 8. & Hệ thống & Hoàn tất ghi dữ liệu và thông báo kết quả thời gian thực \\ \hline
\multirow{4}{*}{\parbox{3cm}{\centering\textbf{Luồng sự kiện}\\\textbf{thay thế}}}
 & \textbf{STT} & \textbf{Thực hiện bởi} & \multicolumn{1}{c|}{\textbf{Hành động}} \\ \cline{2-4}
 & 3a. & Hệ thống & Tệp sai định dạng hoặc thiếu trang tính bắt buộc, báo lỗi và dừng xử lý \\ \cline{2-4}
 & 4a. & Cán bộ Phòng Chính trị & Có dòng lỗi ở bước xem trước, có thể sửa và tải lên tệp đã sửa \\ \cline{2-4}
 & 7a. & Hệ thống & Có lỗi tại bước ghi, huỷ bỏ toàn bộ thay đổi và giữ nguyên dữ liệu cũ \\ \hline
\textbf{Hậu điều kiện} & \multicolumn{3}{m{10.4cm}|}{Các bản ghi hợp lệ trong tệp được ghi vào cơ sở dữ liệu; nhật ký hành động được ghi với tổng kết số dòng thành công và số dòng lỗi; Cán bộ Phòng Chính trị nhận thông báo kết quả thời gian thực.} \\ \hline
\end{tabular}
\caption{Đặc tả use case nhập dữ liệu khen thưởng hàng loạt từ tệp Excel}
\label{tab:uc06}
\end{table}

\begin{table}[H]
\centering
\renewcommand{\arraystretch}{1.1}
\begin{tabular}{|>{\centering\arraybackslash}m{0.8cm}|>{\centering\arraybackslash}m{3.0cm}|>{\raggedright\arraybackslash}m{4.2cm}|>{\centering\arraybackslash}m{1.8cm}|>{\centering\arraybackslash}m{2.8cm}|}
\hline
\textbf{STT} & \textbf{Trường dữ liệu} & \multicolumn{1}{c|}{\textbf{Mô tả}} & \textbf{Bắt buộc} & \textbf{Ví dụ hợp lệ} \\ \hline
1 & Loại khen thưởng & Một trong các loại: danh hiệu hằng năm, khen thưởng đơn vị hằng năm, HCCSVV, HCBVTQ, KNC, HCQKQT, thành tích khoa học & Có & Danh hiệu hằng năm \\ \hline
2 & Tệp Excel & Tệp Excel điền theo mẫu & Có & \texttt{cstdcs-\allowbreak{}2025.xlsx} \\ \hline
\end{tabular}
\caption{Dữ liệu đầu vào cho chức năng nhập Excel hàng loạt}
\label{tab:uc06-input}
\end{table}

\subsection{Đặc tả use case quản lý đơn vị và chức vụ}

\begin{table}[H]
\centering
\renewcommand{\arraystretch}{1.1}
\begin{tabular}{|>{\centering\arraybackslash}m{3.2cm}|>{\centering\arraybackslash}m{0.8cm}|>{\centering\arraybackslash}m{2.8cm}|>{\raggedright\arraybackslash}m{6.2cm}|}
\hline
\textbf{Tên use case} & \multicolumn{3}{m{10.4cm}|}{Quản lý đơn vị (cơ quan đơn vị / đơn vị trực thuộc) và chức vụ} \\ \hline
\textbf{Tác nhân} & \multicolumn{3}{m{10.4cm}|}{Quản trị viên, Cán bộ Phòng Chính trị} \\ \hline
\textbf{Tiền điều kiện} & \multicolumn{3}{m{10.4cm}|}{Đã đăng nhập với vai trò có quyền chỉnh sửa cơ cấu tổ chức.} \\ \hline
\multirow{9}{*}{\parbox{3.2cm}{\centering\textbf{Luồng sự kiện chính}\\\textbf{(Thành công)}}}
 & \textbf{STT} & \textbf{Thực hiện bởi} & \multicolumn{1}{c|}{\textbf{Hành động}} \\ \cline{2-4}
 & 1. & Cán bộ Phòng Chính trị & Thêm cơ quan đơn vị cấp trên: nhập mã đơn vị và tên đơn vị \\ \cline{2-4}
 & 2. & Hệ thống & Kiểm tra mã đơn vị không trùng rồi lưu cơ quan đơn vị \\ \cline{2-4}
 & 3. & Cán bộ Phòng Chính trị & Thêm đơn vị trực thuộc: chọn cơ quan đơn vị cha và nhập thông tin đơn vị \\ \cline{2-4}
 & 4. & Hệ thống & Lưu đơn vị trực thuộc gắn với cơ quan đơn vị cha \\ \cline{2-4}
 & 5. & Cán bộ Phòng Chính trị & Thêm chức vụ cho đơn vị: nhập tên chức vụ, đánh dấu chức vụ quản lý, hệ số chức vụ \\ \cline{2-4}
 & 6. & Hệ thống & Lưu chức vụ kèm ràng buộc duy nhất theo cặp đơn vị và tên chức vụ \\ \cline{2-4}
 & 7. & Cán bộ Phòng Chính trị & Xem cơ cấu tổ chức kèm số quân nhân từng đơn vị \\ \cline{2-4}
 & 8. & Hệ thống & Trả về cơ cấu tổ chức với số quân nhân đã tự động cập nhật \\ \hline
\multirow{4}{*}{\parbox{3cm}{\centering\textbf{Luồng sự kiện}\\\textbf{thay thế}}}
 & \textbf{STT} & \textbf{Thực hiện bởi} & \multicolumn{1}{c|}{\textbf{Hành động}} \\ \cline{2-4}
 & 2a. & Hệ thống & Mã đơn vị đã tồn tại, trả lỗi ``Mã đơn vị đã tồn tại'' \\ \cline{2-4}
 & 4a. & Hệ thống & Cơ quan đơn vị cha không tồn tại, hệ thống báo lỗi \\ \cline{2-4}
 & 8a. & Hệ thống & Cơ quan đơn vị còn đơn vị trực thuộc, chức vụ hoặc quân nhân, hệ thống chặn xoá để bảo toàn dữ liệu \\ \hline
\textbf{Hậu điều kiện} & \multicolumn{3}{m{10.4cm}|}{Cơ cấu tổ chức được cập nhật; số quân nhân từng đơn vị chính xác; quân nhân và chức vụ liên kết đúng đơn vị.} \\ \hline
\end{tabular}
\caption{Đặc tả use case quản lý đơn vị và chức vụ}
\label{tab:uc07}
\end{table}

\subsection{Đặc tả use case sao lưu dữ liệu}

\begin{table}[H]
\centering
\renewcommand{\arraystretch}{1.1}
\begin{tabular}{|>{\centering\arraybackslash}m{3.2cm}|>{\centering\arraybackslash}m{0.8cm}|>{\centering\arraybackslash}m{2.8cm}|>{\raggedright\arraybackslash}m{6.2cm}|}
\hline
\textbf{Tên use case} & \multicolumn{3}{m{10.4cm}|}{Sao lưu dữ liệu định kỳ và thủ công} \\ \hline
\textbf{Tác nhân} & \multicolumn{3}{m{10.4cm}|}{Quản trị viên (qua DevZone), tác vụ tự động theo lịch} \\ \hline
\textbf{Tiền điều kiện} & \multicolumn{3}{m{10.4cm}|}{Tính năng sao lưu đã được bật; Quản trị viên đã xác thực mật khẩu DevZone.} \\ \hline
\multirow{8}{*}{\parbox{3.2cm}{\centering\textbf{Luồng sự kiện chính}\\\textbf{(Thành công)}}}
 & \textbf{STT} & \textbf{Thực hiện bởi} & \multicolumn{1}{c|}{\textbf{Hành động}} \\ \cline{2-4}
 & 1. & Tự động / Quản trị viên & Kích hoạt theo lịch (mặc định 02:00 hằng ngày) hoặc bấm ``Sao lưu ngay'' \\ \cline{2-4}
 & 2. & Hệ thống & Đọc đồng thời 22 bảng gồm dữ liệu nghiệp vụ và nhật ký hệ thống (không gồm bảng thông báo) \\ \cline{2-4}
 & 3. & Hệ thống & Kết xuất dữ liệu từng bảng thành câu lệnh SQL \\ \cline{2-4}
 & 4. & Hệ thống & Ghi tệp sao lưu vào thư mục sao lưu, đặt tên theo thời điểm và loại dữ liệu \\ \cline{2-4}
 & 5. & Hệ thống & Cập nhật mốc sao lưu gần nhất và ghi nhật ký \\ \cline{2-4}
 & 6. & Hệ thống & Xoá các tệp cũ vượt số ngày lưu trữ (mặc định 15) \\ \cline{2-4}
 & 7. & Quản trị viên & Xem trạng thái, đổi lịch hoặc bật/tắt sao lưu trong DevZone \\ \hline
\multirow{4}{*}{\parbox{3cm}{\centering\textbf{Luồng sự kiện}\\\textbf{thay thế}}}
 & \textbf{STT} & \textbf{Thực hiện bởi} & \multicolumn{1}{c|}{\textbf{Hành động}} \\ \cline{2-4}
 & 1a. & Hệ thống & Tính năng sao lưu đang tắt, hệ thống bỏ qua và không tạo tệp \\ \cline{2-4}
 & 4a. & Hệ thống & Lỗi ghi tệp (đầy ổ đĩa, thiếu quyền), hệ thống ghi nhật ký lỗi sao lưu \\ \cline{2-4}
 & 7a. & Hệ thống & Lịch chạy định kỳ không hợp lệ, hệ thống báo lỗi và giữ lịch cũ \\ \hline
\textbf{Hậu điều kiện} & \multicolumn{3}{m{10.4cm}|}{Tệp sao lưu mới được lưu trong thư mục sao lưu; việc khôi phục thực hiện thủ công khi cần.} \\ \hline
\end{tabular}
\caption{Đặc tả use case sao lưu dữ liệu}
\label{tab:uc08}
\end{table}

\subsection{Đặc tả use case quản lý tài khoản}

\begin{table}[H]
\centering
\renewcommand{\arraystretch}{1.1}
\begin{tabular}{|>{\centering\arraybackslash}m{3.2cm}|>{\centering\arraybackslash}m{0.8cm}|>{\centering\arraybackslash}m{2.8cm}|>{\raggedright\arraybackslash}m{6.2cm}|}
\hline
\textbf{Tên use case} & \multicolumn{3}{m{10.4cm}|}{Quản lý tài khoản người dùng} \\ \hline
\textbf{Tác nhân} & \multicolumn{3}{m{10.4cm}|}{Quản trị viên, Cán bộ Phòng Chính trị} \\ \hline
\textbf{Tiền điều kiện} & \multicolumn{3}{m{10.4cm}|}{Đã đăng nhập với vai trò Quản trị viên hoặc Cán bộ Phòng Chính trị.} \\ \hline
\multirow{9}{*}{\parbox{3.2cm}{\centering\textbf{Luồng sự kiện chính}\\\textbf{(Thành công)}}}
 & \textbf{STT} & \textbf{Thực hiện bởi} & \multicolumn{1}{c|}{\textbf{Hành động}} \\ \cline{2-4}
 & 1. & Cán bộ Phòng Chính trị & Tạo tài khoản: nhập tên đăng nhập, chọn vai trò, đơn vị và chức vụ \\ \cline{2-4}
 & 2. & Hệ thống & Băm mật khẩu (dùng mật khẩu mặc định nếu để trống); tạo quân nhân nếu chưa có và lưu tài khoản gắn với quân nhân đó \\ \cline{2-4}
 & 3. & Hệ thống & Thông báo tạo thành công; người dùng đăng nhập lần đầu bằng mật khẩu mặc định \\ \cline{2-4}
 & 4. & Cán bộ Phòng Chính trị & Cập nhật vai trò tài khoản \\ \cline{2-4}
 & 5. & Cán bộ Phòng Chính trị & Đặt lại mật khẩu cho tài khoản \\ \cline{2-4}
 & 6. & Hệ thống & Đặt lại mật khẩu về giá trị mặc định, băm và lưu \\ \cline{2-4}
 & 7. & Cán bộ Phòng Chính trị & Xoá tài khoản (có xác nhận) \\ \cline{2-4}
 & 8. & Hệ thống & Xoá tài khoản; các bản ghi nhật ký liên quan vẫn được giữ lại \\ \hline
\multirow{3}{*}{\parbox{3cm}{\centering\textbf{Luồng sự kiện}\\\textbf{thay thế}}}
 & \textbf{STT} & \textbf{Thực hiện bởi} & \multicolumn{1}{c|}{\textbf{Hành động}} \\ \cline{2-4}
 & 1a. & Hệ thống & Trùng tên đăng nhập, trả lỗi ``Tên đăng nhập đã tồn tại'' \\ \hline
\textbf{Hậu điều kiện} & \multicolumn{3}{m{10.4cm}|}{Tài khoản được tạo/cập nhật/xoá trong cơ sở dữ liệu; người dùng có thể đăng nhập với tài khoản mới.} \\ \hline
\end{tabular}
\caption{Đặc tả use case quản lý tài khoản}
\label{tab:uc09}
\end{table}

\subsection{Đặc tả use case xem nhật ký hệ thống}

\begin{table}[H]
\centering
\renewcommand{\arraystretch}{1.1}
\begin{tabular}{|>{\centering\arraybackslash}m{3.2cm}|>{\centering\arraybackslash}m{0.8cm}|>{\centering\arraybackslash}m{2.8cm}|>{\raggedright\arraybackslash}m{6.2cm}|}
\hline
\textbf{Tên use case} & \multicolumn{3}{m{10.4cm}|}{Xem nhật ký hệ thống} \\ \hline
\textbf{Tác nhân} & \multicolumn{3}{m{10.4cm}|}{Quản trị viên, Cán bộ Phòng Chính trị, Chỉ huy đơn vị} \\ \hline
\textbf{Tiền điều kiện} & \multicolumn{3}{m{10.4cm}|}{Đã đăng nhập với vai trò Quản trị viên, Cán bộ Phòng Chính trị hoặc Chỉ huy đơn vị.} \\ \hline
\multirow{7}{*}{\parbox{3.2cm}{\centering\textbf{Luồng sự kiện chính}\\\textbf{(Thành công)}}}
 & \textbf{STT} & \textbf{Thực hiện bởi} & \multicolumn{1}{c|}{\textbf{Hành động}} \\ \cline{2-4}
 & 1. & Người dùng & Truy cập trang ``Nhật ký hệ thống'' \\ \cline{2-4}
 & 2. & Hệ thống & Hiển thị bảng ghi: thời gian, người thực hiện, vai trò, hành động, nhóm chức năng, mô tả \\ \cline{2-4}
 & 3. & Người dùng & Áp bộ lọc theo nhóm chức năng, hành động, người thực hiện, khoảng thời gian \\ \cline{2-4}
 & 4. & Hệ thống & Lọc theo vai trò: Chỉ huy đơn vị chỉ thấy bản ghi trong phạm vi đơn vị; Cán bộ Phòng Chính trị và Chỉ huy đơn vị không thấy nhật ký sao lưu \\ \cline{2-4}
 & 5. & Hệ thống & Trả về danh sách đã phân trang \\ \cline{2-4}
 & 6. & Người dùng & Đọc mô tả chi tiết của từng bản ghi hiển thị trực tiếp trong bảng \\ \hline
\multirow{3}{*}{\parbox{3cm}{\centering\textbf{Luồng sự kiện}\\\textbf{thay thế}}}
 & \textbf{STT} & \textbf{Thực hiện bởi} & \multicolumn{1}{c|}{\textbf{Hành động}} \\ \cline{2-4}
 & 4a. & Hệ thống & Cán bộ Phòng Chính trị hoặc Chỉ huy đơn vị cố xem nhật ký sao lưu, hệ thống tự lọc bỏ \\ \cline{2-4}
 & 4b. & Hệ thống & Chỉ huy đơn vị chỉ thấy nhật ký do tài khoản trong phạm vi đơn vị họ thực hiện \\ \hline
\textbf{Hậu điều kiện} & \multicolumn{3}{m{10.4cm}|}{Người dùng nắm được lịch sử hành động trong hệ thống, phục vụ kiểm tra và truy vết.} \\ \hline
\end{tabular}
\caption{Đặc tả use case xem nhật ký hệ thống}
\label{tab:uc10}
\end{table}

\section{Yêu cầu phi chức năng}
\label{section:2.4}

Bên cạnh các yêu cầu chức năng, hệ thống phải đáp ứng các yêu cầu phi chức năng nhằm đảm bảo vận hành ổn định, an toàn và phù hợp với đặc thù môi trường quân sự.
Các yêu cầu này được tổ chức theo bốn khía cạnh: hiệu năng và trải nghiệm người dùng, bảo mật, tính nhất quán dữ liệu, khả năng bảo trì và vận hành.

\subsection{Hiệu năng và trải nghiệm người dùng}
\label{subsection:2.4.1}

Hệ thống cần vận hành ổn định, đáp ứng nhanh trong mạng nội bộ Học viện, phục vụ đồng thời cán bộ Phòng Chính trị và chỉ huy các đơn vị, nhất là vào các đợt xét khen thưởng cuối năm.
Các thao tác tra cứu, xử lý thông thường cần phản hồi gần như tức thời. Riêng việc tính lại điều kiện cho toàn bộ quân nhân là tác vụ chạy nền theo lịch định kỳ hoặc khi Quản trị viên chạy thủ công, nên có thể kéo dài hơn mà không ảnh hưởng tới các thao tác khác.

Các danh sách lớn như toàn bộ quân nhân hay khen thưởng theo năm cần được phân trang để tải nhanh.
Những thay đổi như thông báo mới hay thao tác chọn quân nhân vào đề xuất cần phản hồi ngay trên giao diện mà không phải tải lại trang.

\subsection{Bảo mật}
\label{subsection:2.4.2}

Hệ thống lưu dữ liệu cán bộ và sĩ quan trong môi trường quân sự nên bảo mật là yêu cầu ưu tiên hàng đầu.
Cụ thể, hệ thống cần:

\begin{enumerate}
    \item Xác thực người dùng an toàn bằng JWT có cơ chế làm mới luân phiên và thu hồi được từ phía máy chủ; mật khẩu được băm trước khi lưu, không lưu ở dạng văn bản thô.
    \item Phân quyền theo bốn vai trò, giới hạn mỗi vai trò đúng phạm vi chức năng và dữ liệu được phép; riêng Chỉ huy đơn vị chỉ thao tác trong phạm vi đơn vị mình phụ trách.
    \item Kiểm tra mọi dữ liệu đầu vào ở cả phía giao diện lẫn máy chủ, từ chối dữ liệu không hợp lệ kể cả khi yêu cầu được gửi thẳng tới API.
    \item Bảo vệ kết nối khi triển khai, hạn chế các dạng tấn công phổ biến từ phía trình duyệt và chặn việc một địa chỉ gửi quá nhiều yêu cầu trong thời gian ngắn.
    \item Ghi nhật ký đầy đủ mọi thao tác làm thay đổi dữ liệu kèm người thực hiện, thời điểm và nội dung dữ liệu trước và sau khi thay đổi; người thao tác không thể tự xoá dấu vết của mình.
    \item Sao lưu định kỳ và khôi phục được khi cần.
\end{enumerate}

Chi tiết cài đặt các biện pháp này được trình bày tại Chương~\ref{chapter:Methodology} và Chương~\ref{chapter:Experiment}.

\subsection{Tính nhất quán và chính xác của dữ liệu}
\label{subsection:2.4.3}

Một sai lệch nhỏ giữa các bản ghi có thể dẫn tới quyết định sai, chẳng hạn bỏ sót quân nhân đủ điều kiện hoặc đề xuất nhầm chu kỳ.
Vì vậy, tính nhất quán dữ liệu được đặt ngang hàng với bảo mật, với các yêu cầu:

\begin{enumerate}
    \item Các nghiệp vụ ghi trên nhiều bảng phải diễn ra trọn vẹn theo nguyên tắc tất cả hoặc không: khi một bước gặp lỗi, dữ liệu trở về nguyên trạng, không để lại trạng thái dở dang.
    \item Việc tính lại hồ sơ và việc kiểm tra khi phê duyệt phải dùng chung một bộ quy tắc xét điều kiện chuỗi danh hiệu, để hai bên luôn cho cùng một kết luận.
    \item Hồ sơ điều kiện của mỗi quân nhân được suy ra từ dữ liệu gốc, nên phải tự động tính lại mỗi khi dữ liệu gốc thay đổi, để hồ sơ luôn phản ánh đúng dữ liệu hiện tại.
    \item Quan hệ giữa các bảng và các ràng buộc duy nhất phải được giữ toàn vẹn, chẳng hạn mỗi quân nhân chỉ có một bản ghi danh hiệu cho một năm; khi vi phạm, hệ thống báo lỗi rõ ràng thay vì thông báo kỹ thuật khó hiểu.
\end{enumerate}

\subsection{Khả năng bảo trì và vận hành}
\label{subsection:2.4.4}

Hệ thống dự kiến được mở rộng và bảo trì lâu dài nên kiến trúc phải dễ tiếp cận với người tiếp quản.
Mã nguồn được tổ chức phân tầng rõ ràng để khoanh vùng thay đổi và bổ sung loại nghiệp vụ mới mà ít ảnh hưởng tới phần còn lại; mô tả dữ liệu được khai báo tập trung làm nguồn tham chiếu duy nhất; và mọi thao tác làm thay đổi dữ liệu đều có nhật ký truy vết, hỗ trợ rà soát, chẩn đoán khi vận hành.

Về độ sẵn sàng, hệ thống chạy nội bộ trong mạng Học viện, khởi động lại nhanh sau khi cập nhật hoặc gặp sự cố và được sao lưu hằng ngày để giới hạn lượng dữ liệu có thể mất ở mức không quá 24 giờ.

Chương này đã khảo sát hiện trạng công tác khen thưởng tại Học viện, mô hình hoá yêu cầu qua biểu đồ use case tổng quát và các biểu đồ phân rã chức năng, đặc tả chi tiết các use case tiêu biểu và tổng hợp những yêu cầu phi chức năng về hiệu năng, bảo mật và tính nhất quán dữ liệu.
Đây là đầu vào trực tiếp cho việc lựa chọn công nghệ ở chương tiếp theo và thiết kế hệ thống ở các chương sau.

\end{document}
